\documentclass[11pt]{article}
\usepackage[margin=0.8in]{geometry}
\usepackage[T1]{fontenc}
\usepackage{lmodern}
\usepackage{microtype}
\usepackage{parskip}
\usepackage{enumitem}
\usepackage{xcolor}

\definecolor{Ink}{HTML}{1F2937}
\definecolor{Muted}{HTML}{6B7280}

\setlist[itemize]{topsep=3pt,itemsep=3pt,leftmargin=1.2em}

\title{General Field Presentation\\Cliff Notes}
\author{Sahara Kaul \and Kelsey Pattison \and Ahmed Sajid}
\date{CMPUT 414, Winter 2026}

\begin{document}
\color{Ink}
\maketitle

\textbf{Purpose}\\
Quick live-delivery notes. These are for glancing at during the talk, not reading word-for-word.

\textbf{Speaker split}\\
Kelsey: introduction and core challenges\\
Sarah: early generation, VAEs, RAVE, MoVE\\
Ahmed: modern controllable synthesis, diffusion, field trends, conclusion

\section*{Kelsey}

\textbf{Style transfer: what is it?}
\begin{itemize}
\item Start with image analogy: preserve scene structure, change rendering style.
\item In music: preserve content, change timbre, instrumentation, texture, articulation, production character.
\item Mention demo if using one.
\end{itemize}
\textcolor{Muted}{Transition: define content versus style clearly.}

\textbf{Style vs content}
\begin{itemize}
\item Content: melody, rhythm, notes, harmonic path, phrase structure.
\item Style: timbre, instrumentation, texture, dynamics, expressive details.
\item Same content can sound completely different under a new style.
\end{itemize}
\textcolor{Muted}{Transition: now explain why this is harder in music than it sounds.}

\textbf{Why music style transfer is hard}
\begin{itemize}
\item Music is multilevel: melody, harmony, texture, production, room, long-range structure all mixed together.
\item Weak systems create a “coat of paint” result instead of true re-authoring.
\item Genre is broader than timbre: conventions, groove, production norms, listener expectations.
\end{itemize}
\textcolor{Muted}{Transition: one of the key recurring ideas is representation.}

\textbf{Representations}
\begin{itemize}
\item Raw waveforms, spectrograms / log-mel, MIDI, learned latent spaces.
\item Each representation exposes some structure and hides other information.
\item Representation choice matters as much as architecture choice.
\end{itemize}
\textcolor{Muted}{Transition: briefly explain STFT and MIDI as anchor examples.}

\textbf{STFT and MIDI}
\begin{itemize}
\item Fourier Transform: what frequencies exist overall.
\item STFT: what frequencies exist over time, gives a spectrogram.
\item MIDI is easy to manipulate but weak for final timbral realism.
\end{itemize}
\textcolor{Muted}{Transition: this leads to the central field question.}

\textbf{Central field question}
\begin{itemize}
\item How do we change the style or genre of a song while keeping the musical identity recognizable?
\end{itemize}
\textcolor{Muted}{Transition: answer that by naming the field bottlenecks.}

\textbf{Core challenges}
\begin{itemize}
\item Representation problem.
\item Disentanglement problem.
\item Realism problem.
\item Long-form coherence problem.
\item These bottlenecks give a better field structure than just listing models.
\end{itemize}
\textcolor{Muted}{Transition: Sarah picks up with the early audio-generation story.}

\section*{Sarah}

\textbf{WaveNet}
\begin{itemize}
\item Early milestone in realistic neural audio.
\item Autoregressive and slow: one sample at a time.
\item Important historically, but not a practical direct answer to style transfer.
\end{itemize}
\textcolor{Muted}{Transition: this motivates GANs and parallel generation.}

\textbf{Background about GANs}
\begin{itemize}
\item Generator makes fake samples; discriminator spots fake vs real.
\item Often sharper than plain autoencoder outputs.
\item Training is unstable and highly dependent on representation choice.
\end{itemize}
\textcolor{Muted}{Transition: first major audio GAN example is WaveGAN.}

\textbf{WaveGAN}
\begin{itemize}
\item Parallel raw-audio generation, 1D convolutions, phase shuffle.
\item Limits: phase precession and short fixed-length output.
\item Shows raw waveform GANs are possible, but not enough for robust style transfer.
\end{itemize}
\textcolor{Muted}{Transition: GANSynth improves the working space.}

\textbf{GANSynth}
\begin{itemize}
\item Generate log-magnitude + instantaneous-frequency spectrograms instead of raw waveform.
\item More stable than raw phase.
\item Main lesson: representation choice matters.
\end{itemize}
\textcolor{Muted}{Transition: summarize why early generation is only background.}

\textbf{Why early generation is mainly background}
\begin{itemize}
\item WaveNet, WaveGAN, and GANSynth explain how neural audio generation developed.
\item But they do not solve content preservation while changing style.
\item Combined takeaway: generation is possible, representation matters, true transfer remains unsolved.
\end{itemize}
\textcolor{Muted}{Transition: now move into latent structured audio.}

\textbf{VAEs}
\begin{itemize}
\item Useful because they expose a latent space.
\item Problems: posterior collapse, weak bottleneck, sequential complexity.
\item Latent-space idea still matters for style transfer.
\end{itemize}
\textcolor{Muted}{Transition: first talk about RAVE as practical latent audio, then MoVE as more transfer-specific.}

\textbf{RAVE}
\begin{itemize}
\item Not a strong disentanglement model, but makes high-quality latent audio practical.
\item Uses spectral losses, latent sequences, and adversarial sharpening.
\item Best framed as latent audio infrastructure: encode, manipulate, resynthesize.
\end{itemize}
\textcolor{Muted}{Transition: now move to a model that is more directly about transfer.}

\textbf{MoVE}
\begin{itemize}
\item Explicitly about many-to-many musical timbre transfer.
\item Uses VAE latent space, FiLM conditioning, MMD objective.
\item More transfer-specific than RAVE, but still mainly timbre transfer rather than full genre remastering.
\end{itemize}
\textcolor{Muted}{Transition: this latent section matters because it gets much closer to actual style transfer than the older GAN section.}

\section*{Ahmed}

\textbf{Modern controllable synthesis}
\begin{itemize}
\item New question: once we have a usable representation and some notion of style, how do we render convincing controllable audio?
\end{itemize}
\textcolor{Muted}{Transition: start with BigVGAN and FiLM because they do different jobs.}

\textbf{BigVGAN}
\begin{itemize}
\item Universal neural vocoder.
\item Converts intermediate representation to waveform audio.
\item Role in the story: realism stage, not transfer logic.
\end{itemize}
\textcolor{Muted}{Transition: FiLM is about conditioning, not waveform realism.}

\textbf{FiLM}
\begin{itemize}
\item Feature-wise Linear Modulation.
\item Style vector scales and shifts hidden features.
\item Good way to inject style information throughout the network.
\end{itemize}
\textcolor{Muted}{Transition: now move into diffusion itself.}

\textbf{DDPM}
\begin{itemize}
\item Forward noising process, learned reverse denoising process.
\item More stable than many adversarial setups.
\item Important because gradual refinement gives more control.
\end{itemize}
\textcolor{Muted}{Transition: classifier-free guidance and SDEdit make that control more explicit.}

\textbf{Classifier-Free Guidance}
\begin{itemize}
\item Guidance scale \(w\) controls how strongly generation follows the target condition.
\item Lower guidance preserves naturalness; higher guidance pushes harder toward style.
\end{itemize}
\textcolor{Muted}{Transition: SDEdit is the editing version of the same idea.}

\textbf{SDEdit}
\begin{itemize}
\item Start from a noised source instead of pure noise.
\item Lets the model preserve the source skeleton while still changing the sound world.
\end{itemize}
\textcolor{Muted}{Transition: this sets up AudioLDM.}

\textbf{AudioLDM}
\begin{itemize}
\item Latent diffusion for audio.
\item Important because diffusion happens in a learned audio latent space instead of raw waveform space.
\item Strong modern direction for controllable audio transformation and re-authoring.
\end{itemize}
\textcolor{Muted}{Transition: now summarize why diffusion is the modern direction.}

\textbf{Why diffusion is the modern generative direction}
\begin{itemize}
\item Progressive generation.
\item Strong conditioning.
\item Explicit guidance control.
\item Source anchoring through SDEdit-style editing.
\end{itemize}
\textcolor{Muted}{Transition: end by summarizing the overall field movement.}

\textbf{Overview of field trends}
\begin{itemize}
\item More useful arc: disentanglement, controllable latent audio, diffusion-based re-authoring, long-form coherence.
\item Field moved from raw generation toward structured representations and explicit control.
\end{itemize}
\textcolor{Muted}{Transition: final lesson.}

\textbf{Conclusion}
\begin{itemize}
\item Music style transfer is not just an audio-generation problem.
\item It is also a structure-preservation, representation, conditioning, synthesis, and coherence problem.
\item Clean ending: the field is moving from surface transfer toward structure-aware genre reconstruction.
\end{itemize}

\end{document}
